\title{DeepSeekMath: Pushing the Limits of Mathematical Reasoning in Open Language Models}

\begin{abstract}
Mathematical reasoning represents a formidable challenge for language models, stemming from its inherent complexity and structured demands. In this study, we present DeepSeekMath~7B, a model that extends the pre-training of DeepSeek-Coder-Base-v1.5 7B by leveraging 120B math-specific tokens drawn from Common Crawl, in conjunction with both natural language and code corpora. DeepSeekMath~7B demonstrates a notable score of 51.7\% on the competitive MATH benchmark, attaining this result independently of external toolkits or voting mechanisms—narrowing the gap with leading models such as Gemini-Ultra and GPT-4. Further, when evaluated with self-consistency using 64 samples, DeepSeekMath~7B achieves a score of 60.9\% on MATH. This enhanced mathematical reasoning is ascribed to two principal innovations: firstly, we exploit the richness of open web data through a carefully constructed data selection process; secondly, we propose Group Relative Policy Optimization (GRPO), a new variant of Proximal Policy Optimization (PPO) that both bolsters mathematical reasoning and improves PPO’s memory efficiency.
\end{abstract}

\section{Introduction}

The advent of large language models (LLM) has fundamentally transformed mathematical reasoning within artificial intelligence, enabling rapid progress on quantitative reasoning benchmarks~\citep{MATH} as well as geometry-focused evaluations~\citep{trinh2024solving}. These LLMs have become invaluable resources for aiding human efforts in tackling challenging mathematical tasks~\citep{tao}. Nonetheless, state-of-the-art solutions such as GPT-4~\citep{gpt4} and Gemini-Ultra~\citep{gemini} remain closed-source, resulting in a notable disparity, as open-source alternatives still fall significantly short in terms of performance.

This work presents DeepSeekMath, a specialized language model designed for mathematical reasoning, which exhibits mathematical performance well beyond that of existing open-source models and nears the accuracy of GPT-4 across standardized academic evaluations. Central to this achievement is the construction of the DeepSeekMath Corpus, a vast, meticulously curated pre-training dataset containing 120B math tokens. The dataset is collected from Common Crawl (CC) via a fastText-based classifier \citep{joulin2016fasttext}. Initially, this classifier is trained on positive samples from OpenWebMath \citep{openwebmath} and a broad array of unrelated web documents serving as negative samples. The trained classifier is subsequently deployed to extract further math-relevant content from CC, which is subsequently validated and refined through manual annotation. This process iterates, with the improved, annotated data used to retrain the classifier, enhancing its discriminative capabilities. Empirical results demonstrate the effectiveness and quality of this large-scale corpus: our primary model, DeepSeekMath-Base 7B, obtains 64.2\% accuracy on GSM8K \citep{gsm8k} and 36.2\% on the advanced MATH benchmark \citep{MATH}, thereby surpassing the performance of Minerva 540B \citep{minerva}. The DeepSeekMath Corpus also encompasses multiple languages, yielding noticeable gains on Chinese mathematical evaluations \citep{wei2023cmath,agieval}. We view our methodology for curating mathematical datasets as a preliminary yet valuable contribution to the field, anticipating substantial opportunities for advancement moving forward.

DeepSeekMath-Base is built upon DeepSeek-Coder-Base-v1.5 7B \citep{deepseek-coder}, as empirical observations suggest that initializing from a model pre-trained on code yields superior performance compared to initializing from a general-purpose language model. Moreover, we find that mathematical pre-training enhances performance not only on math-specific tasks, but also on general reasoning benchmarks such as MMLU \citep{mmlu} and BBH \citep{bbh}, signifying improvements that extend beyond mathematics to broader reasoning abilities.

Following pre-training, we subject DeepSeekMath-Base to a mathematical instruction tuning process, incorporating chain-of-thought \citep{cot}, program-of-thought \citep{pot,pal}, and tool-augmented reasoning datasets \citep{tora}. The resulting model, DeepSeekMath-Instruct 7B, consistently surpasses other 7B parameter models and attains performance levels comparable to open-source instruction-tuned models at the 70B scale.

Additionally, we present Group Relative Policy Optimization (GRPO), an alternative reinforcement learning (RL) method inspired by Proximal Policy Optimization (PPO) \citep{schulman2017proximal}. GRPO dispenses with the critic model and instead derives baselines from group-based scoring, thereby considerably lowering the resource demands of training. Utilizing only a selective portion of English instruction tuning data, GRPO markedly elevates the strong baseline set by DeepSeekMath-Instruct, achieving substantial gains both on in-domain datasets (GSM8K: 82.9\% $\rightarrow$ 88.2\%, MATH: 46.8\% $\rightarrow$ 51.7\%) and on out-of-domain tasks (e.g., CMATH: 84.6\% $\rightarrow$ 88.8\%) during RL fine-tuning. We also propose a unifying framework to interpret a range of RL methods, including Rejection Sampling Fine-Tuning (RFT) \citep{yuan2023scaling}, Direct Preference Optimization (DPO) \citep{dpo}, PPO, and GRPO. Under this unified perspective, we observe that all aforementioned techniques are variants or simplified forms of RL. Our comprehensive experimental study includes analyses such as comparing online and offline training, evaluating outcome versus process supervision, and assessing single-step against iterative RL approaches to scrutinize the critical facets of this unifying framework. Finally, we provide an explanation of why RL substantially benefits instruction-tuned models, and outline future avenues for optimizing RL efficacy within this conceptual paradigm.

\subsection{Contributions}
This work presents advances in large-scale mathematical pre-training, as well as an investigation into reinforcement learning approaches.

\textbf{Math Pre-Training at Scale}

\begin{itemize}[topsep=0pt]
    \item We demonstrate that the openly available Common Crawl resource contains substantial mathematical content. Leveraging a rigorously constructed data selection process, we developed the DeepSeekMath Corpus—a high-quality collection comprising 120B tokens drawn from mathematics-rich web pages. This corpus is approximately 7 times larger than the dataset utilized by Minerva \citep{minerva}, and nearly 9 times greater than the size of the newly published OpenWebMath dataset \citep{openwebmath}.

    \item Our DeepSeekMath-Base 7B model, after pre-training, delivers performance on par with the substantially larger Minerva 540B model \citep{minerva}. This outcome suggests that scaling up model parameters is not the sole determinant of mathematical reasoning abilities; rather, a smaller model trained on a curated, high-quality mathematical dataset can also excel.

    \item We report experimental results on the effects of sequential training: specifically, models that undergo code pre-training prior to mathematical training exhibit enhanced proficiency in tackling mathematical problems, regardless of the use of external tools. These results address the persisting question: \textit{does code training enhance reasoning skills?} Our findings affirm that it does—particularly in the context of mathematical reasoning.

    \item While the inclusion of arXiv papers is a standard practice for training, especially in mathematical domains, we observe that this addition yields no significant gains across the mathematical benchmarks considered in this study.
\end{itemize}

\textbf{Exploration and Analysis of Reinforcement Learning}

\begin{itemize}[topsep=0pt]
    \item We present Group Relative Policy Optimization (GRPO), a reinforcement learning algorithm characterized by both efficiency and efficacy. GRPO eliminates the need for a critic network by deriving the baseline directly from group scores, thereby considerably cutting down on computational requirements relative to Proximal Policy Optimization (PPO).
    \item We provide evidence that GRPO brings marked improvements to our instruction-tuned model, DeepSeekMath-Instruct, leveraging only instruction-tuning data. Additionally, we note that reinforcement learning with GRPO contributes to enhanced generalization to out-of-domain tasks.
    \item We offer an integrated conceptual framework that encompasses a range of algorithms, including RFT, DPO, PPO, and GRPO. Our empirical analysis covers a variety of scenarios—such as the comparison between online and offline learning, outcome-based versus process-oriented supervision, and single-turn versus iterative reinforcement learning—enabling a thorough exploration of the critical aspects of this unified approach.
    \item Building on the established unified framework, we investigate the underlying causes behind the success of reinforcement learning strategies and outline several promising research avenues for achieving more effective reinforcement learning in large language models (LLMs).
\end{itemize}

\subsection{Summary of Evaluations and Metrics}

\begin{itemize}[topsep=0pt]
    \item \textbf{English and Chinese Mathematical Reasoning}: We undertake a thorough evaluation of our models on both English and Chinese mathematical benchmarks, spanning questions from elementary to collegiate difficulty. The English benchmarks utilized include GSM8K \citep{gsm8k}, MATH \citep{MATH}, SAT \citep{llemma}, OCW Courses \citep{minerva}, and MMLU-STEM \citep{mmlu}; for Chinese, we employ MGSM-zh \citep{mgsm}, CMATH \citep{wei2023cmath}, Gaokao-MathCloze \citep{agieval}, and Gaokao-MathQA \citep{agieval}. Our evaluation covers both the models’ capacity to independently generate text-based solutions without the aid of external tools and their competence in solving tasks programmatically using Python.

    On the suite of English tasks, DeepSeekMath-Base achieves results comparable to the proprietary Minerva 540B \citep{minerva} and decisively outperforms all currently available open-source models—including Mistral 7B \citep{mistral} and Llemma-34B \citep{llemma}—regardless of whether those models were subjected to mathematics-specific pre-training, with differences frequently being substantial. Impressively, DeepSeekMath-Base achieves even greater success on the Chinese benchmarks, which may be attributed to our strategy of integrating high-quality, multilingual pre-training data rather than restricting to English as seen in previous studies \citep{minerva,llemma}. When further enhanced with mathematical instruction tuning and reinforcement learning, DeepSeekMath-Instruct and DeepSeekMath-RL exhibit robust capabilities, achieving an unprecedented open-source benchmark: surpassing 50\% accuracy on the competition-grade MATH dataset.

\item \textbf{Formal Mathematics}: The evaluation of DeepSeekMath-Base is performed using the informal-to-formal theorem proving benchmark introduced by \citep{dsp_proof}, leveraging miniF2F \citep{minif2f} where Isabelle \citep{isabelle} serves as the designated proof assistant. In this context, DeepSeekMath-Base achieves notable performance in few-shot autoformalization.

\item \textbf{Natural Language Understanding, Reasoning, and Code}: To thoroughly assess the model’s proficiency in broad language understanding, reasoning, and programming, we examine DeepSeekMath-Base on the Massive Multitask Language Understanding (MMLU) suite \citep{mmlu}, which spans 57 diverse multiple-choice tasks, as well as the BIG-Bench Hard (BBH) collection \citep{bbh}, comprising 23 tasks demanding primarily complex, multi-step reasoning. The evaluation extends to coding benchmarks HumanEval \citep{codex} and MBPP \citep{mbpp}, which are prevalent in code model assessments. The results indicate that mathematical pre-training enhances both comprehension and reasoning capabilities in these domains.
\end{itemize}

\section{Math Pre-Training}

\subsection{Data Collection and Decontamination} This section details the methodology adopted to assemble the DeepSeekMath Corpus utilizing data from Common Crawl. As shown in Figure \ref{fig:data_collect}, we employ an iterative data collection pipeline designed to systematically amass a substantial mathematical text corpus from Common Crawl, initially seeded with a small, curated set of high-quality mathematics data. Notably, this pipeline can be generalized to extract data for other fields, such as programming.

Our process begins by designating OpenWebMath \citep{openwebmath}, a repository of reputable mathematical web texts, as the seed set. Leveraging this corpus, we train a fastText classifier \citep{joulin2016fasttext} to retrieve additional mathematical web pages resembling OpenWebMath from Common Crawl. For the training phase, we randomly sample 500,000 items from the seed corpus as positive instances and an equivalent number of randomly selected Common Crawl web pages as negative examples. Training is executed with an open-source fastText implementation,\footnote{\url{https://fasttext.cc}} specifying a vector size of 256, a learning rate of 0.1, a maximum word n-gram length of 3, a minimum frequency threshold of 3, and 3 epochs. To manage data scale, we apply URL-based deduplication and near-deduplication on the initial Common Crawl snapshot, reducing it to 40B HTML documents. Using the trained fastText model, we identify and retrieve mathematical content from this deduplicated set. To exclude subpar materials, web pages are scored via the fastText classifier, and only those with the highest scores are retained. The final volume of curated content is determined by pre-training experiments at several thresholds—top 40B, 80B, 120B, and 160B tokens—with the first round opting for the top 40B tokens to be preserved.

Following the initial round of data collection, a significant portion of mathematical web pages remains undiscovered, primarily due to the limited diversity within the set of positive examples used to train the fastText model. To address this, we augment the seed dataset by incorporating new mathematical web sources, thereby enhancing the effectiveness of the fastText model. Concretely, the Common Crawl dataset is partitioned into non-overlapping domains, each defined by a shared base URL. For every domain, we compute the proportion of web pages obtained in the initial collection phase. Domains where more than 10\% of web pages were gathered are designated as math-related (for example, \url{mathoverflow.net}). Subsequently, we manually review and annotate those URLs within these domains that correspond to mathematical content (such as \url{mathoverflow.net/questions}). Any uncollected web pages under these URLs are then added to our seed set. This iterative enrichment yields a broader set of positive examples, allowing the retrained fastText model to recover more mathematical data in later rounds. Ultimately, after four such iterations, our corpus comprises 35.5 million mathematical web pages, amounting to a total of 120 billion tokens. Analysis after the fourth iteration reveals that about 98\% of the data had already been acquired in the third round, prompting the decision to stop further data collection.

To prevent contamination of downstream benchmarks, we adopt the methodology from \cite{deepseek-coder} to filter out web pages that include questions or answers found in English mathematical benchmarks such as GSM8K \citep{gsm8k} and MATH \citep{MATH}, as well as Chinese benchmarks like CMATH \citep{wei2023cmath} and AGIEval \citep{agieval}. The applied filtering rule is as follows: any text fragment that contains a 10-gram substring matching exactly with any segment from these benchmarks is excised from the math training corpus. For benchmark content shorter than 10-grams but at least 3-grams in length, we similarly apply exact substring matching to exclude potentially contaminated web pages.

\subsection{Evaluating the Quality of the DeepSeekMath~Corpus}

To assess the relative quality of the DeepSeekMath~Corpus, we conduct pre-training experiments comparing it to several contemporary mathematical corpora:

\begin{itemize}[topsep=0pt]
    \item \textbf{MathPile} \citep{mathpile}: a composite corpus (8.9B tokens) comprising material from textbooks, Wikipedia, ProofWiki, CommonCrawl, StackExchange, and arXiv—with arXiv documents contributing over 85\% of the total;
    \item \textbf{OpenWebMath} \citep{openwebmath}: a dataset of 13.6B tokens of mathematical content extracted from CommonCrawl;
    \item \textbf{Proof-Pile-2} \citep{llemma}: this collection merges OpenWebMath, AlgebraicStack (10.3B tokens of math-related code), and arXiv articles (28.0B tokens). In our experiments with Proof-Pile-2, we follow the data mixing ratios (arXiv:Web:Code = 2:4:1) reported in \cite{llemma}.
\end{itemize}

\subsubsection{Training Setting}
\label{sec:quality-policy}
We perform mathematical training on a foundational language model possessing 1.3B parameters, architecturally aligned with the DeepSeek LLMs \citep{deepseek-llm}, referred to here as DeepSeek-LLM 1.3B. Individual models are trained separately on each distinct mathematical dataset for a total of 150B tokens. All training runs utilize the streamlined and efficient HAI-LLM \citep{haillm} training system. In accordance with DeepSeek LLMs’ established methodology, we employ the AdamW optimizer \citep{adamW} with hyperparameters set to $\beta_1=0.9$, $\beta_2=0.95$, and $\mathrm{weight\_decay}=0.1$. The learning rate follows a staged schedule: it ramps up over the first 2,000 steps (warmup), decays to 31.6\% of the maximum by the time 80\% of training has elapsed, and further decays to 10.0\% of the maximum after 90\% of total training. The upper bound for the learning rate is fixed at 5.3e-4, with training conducted in batches of 4 million tokens and a sequence length of 4K.

\subsubsection{Evaluation Results}

\textbf{The DeepSeekMath~Corpus distinguishes itself with high quality, comprehensive multilingual content, and unparalleled scale.}

\begin{itemize}[topsep=0pt]
    \item \textbf{High-quality}:
    We assess downstream mathematical ability across 8 benchmark datasets, employing few-shot chain-of-thought prompting \cite{cot}. As illustrated in Table \ref{tab:corpora_comparison}, the model trained on the DeepSeekMath~Corpus clearly outperforms others. Figure \ref{fig:corpora_comparisons} further indicates that after 50B tokens of training (which constitutes one epoch over Proof-Pile-2), DeepSeekMath-trained models exhibit superior results relative to those trained on Proof-Pile-2, reflecting the higher average quality of DeepSeekMath~Corpus content.
    \item \textbf{Multilingual}:
    The DeepSeekMath~Corpus incorporates materials in several languages, with English and Chinese representing the majority. Table \ref{tab:corpora_comparison} demonstrates that exposure to DeepSeekMath~Corpus data significantly boosts mathematical reasoning proficiency in both English and Chinese. In contrast, most current mathematical corpora, which heavily favor English, afford limited gains—and can even impair performance—when it comes to Chinese-language mathematical reasoning.
    \item \textbf{Large-scale}:
    The size of the DeepSeekMath~Corpus greatly surpasses that of other mathematical datasets. As seen in Figure \ref{fig:corpora_comparisons}, DeepSeek-LLM 1.3B, after being trained on DeepSeekMath~Corpus, experiences a sharper increase and more persistent improvements in performance. Conversely, the baseline corpora are markedly smaller, have undergone numerous re-trainings, and as a result, models trained on them quickly reach a point of stagnation in their learning progress.
\end{itemize}

\subsection{Training and Evaluating DeepSeekMath-Base 7B}

Here, we present DeepSeekMath-Base 7B, a foundational model distinguished by its advanced reasoning skills, particularly within the domain of mathematics. The initialization of our model uses DeepSeek-Coder-Base-v1.5 7B \citep{deepseek-coder}, after which it undergoes training on a dataset comprising 500B tokens. The dataset composition is detailed as follows: DeepSeekMath Corpus constitutes 56\%, AlgebraicStack provides 4\%, arXiv accounts for 10\%, Github code makes up 20\%, and the remaining 10\% is drawn from Common Crawl natural language data in English and Chinese. For training, we adhere largely to the configuration in Section \ref{sec:quality-policy}, with modifications to set the peak learning rate at 4.2e-4 and to use a batch size of 10M tokens.

A thorough evaluation is conducted to measure the mathematical proficiency of DeepSeekMath-Base 7B. This includes its capacity to generate fully self-contained mathematical solutions independently of external resources, tackle mathematical problems with the aid of tools, and carry out formal theorem proving. Additionally, we explore the broader capabilities of the base model, such as its proficiency in natural language understanding, logical reasoning, and programming.

\paragraph{Mathematical Problem Solving with Step-by-Step Reasoning}

We assess the problem-solving performance of DeepSeekMath-Base through few-shot chain-of-thought prompting \citep{cot} across eight benchmarks available in both English and Chinese. The evaluation spans tasks involving quantitative reasoning—such as those found in GSM8K \citep{gsm8k}, MATH \citep{MATH}, and CMATH \citep{wei2023cmath}—as well as multiple-choice questions, exemplified by benchmarks like MMLU-STEM \citep{mmlu} and Gaokao-MathQA \citep{agieval}. Together, these benchmarks represent a broad spectrum of mathematical domains, ranging from elementary mathematics to advanced topics encountered at the college level.

As indicated in Table \ref{tab:base_math_cot}, DeepSeekMath-Base 7B establishes new state-of-the-art results across all eight evaluated benchmarks within the category of open-source base models. This includes both the highly popular Mistral 7B model \citep{mistral} and the recently introduced Llemma 34B \citep{llemma}, which was trained on the Proof-Pile-2 dataset \citep{llemma} for mathematical reasoning. On the challenging MATH dataset, which reflects competition-level difficulty, DeepSeekMath-Base achieves a performance lead exceeding 10\% absolute compared to all prior open-source base models. Moreover, it surpasses Minerva 540B \citep{minerva}—a proprietary, large-scale base model constructed upon PaLM \citep{palm} and further specialized with mathematical corpora—even though Minerva contains 77 times more parameters.

\paragraph{Mathematical Problem Solving with Tool Use} For assessing program-supported mathematical reasoning, we utilize the GSM8K and MATH datasets, employing the few-shot program-of-thought prompting approach \citep{pot,pal}. Here, models are instructed to resolve each problem by producing Python scripts that may employ libraries like \textit{math} and \textit{sympy} to address sophisticated calculations. The final program output is then treated as the solution. As evidenced in Table \ref{tab:base_math_pot_proof}, DeepSeekMath-Base 7B surpasses the previously leading open-source model, Llemma 34B, in this context.

\paragraph{Formal Mathematics}

The automation of formal proofs has become increasingly recognized for enhancing both the reliability and productivity of mathematical validation. We test DeepSeekMath-Base 7B on the informal-to-formal proof generation task introduced by \citep{dsp_proof}, where the objective is to translate an informal conjecture, its formal statement, and a human-readable proof into a formalized proof. For this, we use the miniF2F benchmark \citep{minif2f}, which focuses on Olympiad-style mathematical formalization tasks. Formal proofs in Isabelle are generated for each instance via few-shot prompting. Following the methodology in \cite{dsp_proof}, we generate initial proof outlines with the model and employ the automated prover Sledgehammer \citep{sledgehammer} to complete any missing argumentation. Table \ref{tab:base_math_pot_proof} demonstrates the robust performance of DeepSeekMath-Base 7B in automatic proof formalization.

\paragraph{Natural Language Understanding, Reasoning, and Code}

We further assess the model's abilities in natural language understanding using the MMLU benchmark \citep{mmlu}, logical reasoning using BBH \citep{bbh}, and programming proficiency with HumanEval \citep{codex} and MBPP \citep{mbpp}. As seen in Table \ref{tab:understanding_reasoning_code}, DeepSeekMath-Base 7B registers pronounced improvements over its predecessor, DeepSeek-Coder-Base-v1.5 \citep{deepseek-coder}, particularly for the MMLU and BBH datasets, highlighting the advantages of targeted math pretraining for general language comprehension and reasoning. Furthermore, by including code-oriented data during continued training, DeepSeekMath-Base 7B retains parity with DeepSeek-Coder-Base-v1.5 on programming benchmarks. Overall, DeepSeekMath-Base 7B substantially outperforms the generalist model Mistral 7B \citep{mistral} across all three benchmarks addressing reasoning and coding skills.

\section{Supervised Fine-Tuning}

\subsection{SFT Data Curation}

We assemble a dataset for mathematical instruction-tuning that encompasses both English and Chinese mathematical questions. These problems span a range of topics and degrees of difficulty, and each question is provided with an answer formatted using chain-of-thought (CoT) \citep{cot}, program-of-thought (PoT) \citep{pot,pal}, or tool-augmented reasoning methodologies \citep{tora}. The curated training corpus comprises 776,000 instances.

\begin{itemize}[topsep=0pt]
    \item \textbf{English mathematical datasets}: We supply tool-augmented answers to items from GSM8K and MATH, and additionally incorporate selected MathInstruct \citep{MathInstruct} problems as well as those from the Lila-OOD \citep{lila} training split, each annotated with either CoT or PoT solutions. This set covers an array of mathematical domains such as algebra, calculus, number theory, geometry, and probability.
    \item \textbf{Chinese mathematical datasets}: We aggregate K-12 mathematics problems written in Chinese, representing 76 distinct subfields, including topics like linear equations. Corresponding solutions are supplied using both CoT and tool-augmented approaches.
\end{itemize}

\subsection{Training and Evaluating DeepSeekMath-Instruct 7B}

This section details the development of DeepSeekMath-Instruct 7B, a model refined through instruction tuning using DeepSeekMath-Base as its foundation. During training, we concatenate examples at random to fill up to a 4K-token context window. The optimization is conducted for 500 training steps, with a batch size set at 256 and a fixed learning rate of $5 \times 10^{-5}$.

We assess mathematical capabilities both with and without external tool support across four benchmarks in English and Chinese designed for quantitative reasoning. For comparative evaluation, our system is positioned alongside prominent contemporary models:

\begin{itemize}[topsep=0pt]
    \item \textbf{Closed-source models} evaluated include: (1) members of the GPT lineup, in particular, GPT-4 \citep{gpt4} and GPT-4 Code Interpreter~\footnote{\url{https://openai.com/blog/chatgpt-plugins\#\#code-interpreter}} which exhibit top-tier performance, (2) Gemini Ultra and Pro \citep{gemini}, (3) Inflection-2 \citep{inflection-2}, (4) Grok-1~\footnote{\url{https://x.ai/model-card}}, as well as the newest systems developed by Chinese firms such as (5) Baichuan-3~\footnote{\url{https://www.baichuan-ai.com}}, and (6) the advanced GLM-4~\footnote{\url{https://open.bigmodel.cn/dev/api\#glm-4}}
    representing the GLM series \citep{glm}.
\end{itemize}

The aforementioned models are primarily designed for general use, with most having been subjected to various alignment techniques.

\item \textbf{Open-source models} comprise: standard models such as (1) DeepSeek-LLM-Chat 67B \citep{deepseek-llm}, (2) Qwen 72B \citep{qwen}, (3) SeaLLM-v2 7B \citep{seallm}, and (4) ChatGLM3 6B \citep{chatglm3}. In addition, several models feature targeted enhancements in mathematics, including (5) InternLM2-Math 20B~\footnote{\url{https://github.com/InternLM/InternLM-Math}}, which extends InternLM2 with supplementary mathematics training and subsequent instruction tuning, (6) Math-Shepherd-Mistral 7B, which integrates PPO training \citep{schulman2017proximal} with Mistral 7B \citep{mistral} utilizing a reward model guided by process supervision, (7) the WizardMath collection \citep{wizardmath}, aimed at strengthening the mathematical capabilities of Mistral 7B and Llama-2 70B \citep{llama2} via evolve-instruct (an AI-generated instruction tuning approach) and PPO, predominantly using problems from GSM8K and MATH datasets, (8) MetaMath 70B \citep{metamath}, which is based on Llama-2 70B but further tuned on a broadened version of GSM8K and MATH, (9) ToRA 34B \cite{tora}, derived from CodeLlama 34B and optimized for tool-augmented mathematical reasoning, and (10) MAmmoTH 70B \citep{MathInstruct}, developed by instruction tuning Llama-2 70B with MathInstruct data.
\end{itemize}

According to Table \ref{tab:sft_rl_math}, when tool usage is not permitted during evaluation, DeepSeekMath-Instruct 7B stands out in step-by-step reasoning ability. Remarkably, on the competition-tier MATH dataset, this model outperforms every other open-source solution as well as most closed-source systems (such as Inflection-2 and Gemini Pro), achieving an absolute margin of over 9\%. This includes models that possess considerably greater scale (for example, Qwen 72B) or that have benefitted from mathematics-specific reinforcement learning techniques (such as WizardMath-v1.1 7B). DeepSeekMath-Instruct demonstrates performance comparable to leading proprietary Chinese systems GLM-4 and Baichuan-3 on the MATH dataset, yet remains behind the results achieved by GPT-4 and Gemini Ultra.

When evaluation conditions permit both natural language reasoning and programmatic tool use for solving problems, DeepSeekMath-Instruct 7B achieves accuracy rates near 60\% on the MATH dataset, surpassing all publicly available models. In assessments on other benchmarks, the model maintains performance on par with DeepSeek-LLM-Chat 67B, the preceding open-source leader despite having only a tenth of the parameter count.

\section{Reinforcement Learning}

\subsection{Group Relative Policy Optimization}

The application of reinforcement learning (RL) has been shown to further elevate the mathematical reasoning skills of large language models (LLMs) following Supervised Fine-Tuning (SFT) \citep{wang2023math,wizardmath}. Here, we describe our streamlined and high-performing RL method: Group Relative Policy Optimization (GRPO).

\subsubsection{From PPO to GRPO} Proximal Policy Optimization (PPO) \citep{schulman2017proximal} is a prominent actor-critic RL approach often leveraged during the RL fine-tuning phase for LLMs \citep{ouyang2022training}. It seeks to optimize model behavior by maximizing the following surrogate objective function:

\begin{equation}
\footnotesize
    \mathcal{J}_{PPO}(\theta) = \mathbb{E}{[q \sim P(Q), o \sim \pi_{\theta_{old}}(O|q)]} \frac{1}{|o|} \sum_{t=1}^{|o|} \min \left[ \frac{\pi_\theta(o_{t} | q, o_{<t})}{\pi_{\theta_{old}}(o_{t} | q, o_{<t})} A_{t}, \text{clip} \left( \frac{\pi_\theta(o_{t} | q, o_{<t})}{\pi_{\theta_{old}}(o_{t} | q, o_{<t})},

In this context, $\pi_{\theta}$ denotes the current policy model, while $\pi_{\theta_{old}}$ represents the policy model prior to updates. The variables $q$ and $o$ refer, respectively, to questions drawn from the question dataset and outputs sampled under the old policy $\pi_{\theta_{old}}$. The hyper-parameter $\varepsilon$ is associated with the clipping mechanism in PPO, designed to help stabilize the learning process. The advantage $A_t$ is computed using Generalized Advantage Estimation (GAE) \citep{gae}, leveraging the sequence of rewards $\{r_{\ge t}\}$ along with predictions from a learned value function $V_{\psi}$. As a result, PPO necessitates the concurrent training of a value function alongside the policy network. To prevent the reward model from being overly optimized, a typical remedy is to include a KL penalty at each token, comparing the current policy against a reference model in the reward computation \citep{ouyang2022training}, formalized as

\begin{equation}
    r_{t} = r_\varphi(q, o_{\le t}) - \beta \log\frac{\pi_{\theta}(o_{t}|q, o_{<t})}{\pi_{ref}(o_{t}|q, o_{<t})},
\label{eq:PPO-reward}
\end{equation}

where $r_\varphi$ is the output from the reward model, $\pi_{ref}$ is a reference policy (often the original SFT model), and $\beta$ modulates the KL penalty's influence.

Since PPO’s value function is generally an additional model of similar scale to the policy model, it introduces significant demands on memory and computation. During RL optimization, the value function operates as a baseline in the advantage computation for reducing variance. However, in the case of LLMs, reward models typically assign a score only to the final token, making it challenging to train a token-level accurate value function. To overcome these issues, as depicted in Figure \ref{fig:grpo}, we introduce Group Relative Policy Optimization (GRPO). Unlike PPO, GRPO removes the dependency on auxiliary value function estimation and, instead, defines the baseline as the mean reward across multiple outputs sampled for a single question. Concretely, for each question $q$, GRPO draws a batch of outputs $\{o_1, o_2, \cdots, o_G\}$ from $\pi_{\theta_{old}}$ and optimizes the new policy by maximizing the following objective:

\begin{equation}
\footnotesize
\begin{split}
    \mathcal{J}_{GRPO}(\theta) &= \mathbb{E}{[q \sim P(Q), \{o_i\}_{i=1}^G \sim \pi_{\theta_{old}}(O|q)]}  \\
    & \frac{1}{G}\sum_{i=1}^G\frac{1}{|o_i|} \sum_{t=1}^{|o_i|} \left\{ \min \left[ \frac{\pi_\theta(o_{i,t} | q, o_{i,<t})}{\pi_{\theta_{old}}(o_{i,t} | q, o_{i,<t})} \hat{A}_{i,t}, \text{clip} \left( \frac{\pi_\theta(o_{i,t} | q, o_{i,<t})}{\pi_{\theta_{old}}(o_{i,t} | q, o_{i,<t})}, 1 - \varepsilon, 1 + \varepsilon \right)  \hat{A}_{i,t} \right] - \beta \mathbb{D}_{KL}\left[\pi_{\theta} || \pi_{ref}\right]\right\} ,
\end{split}
\label{eq:GRPO-obj}
\end{equation}

Here, $\varepsilon$ and $\beta$ function as hyper-parameters, while $\hat{A}_{i,t}$ denotes the advantage computed by considering only the relative rewards among the outputs in the same group—a method described in more detail in subsequent sections. This grouping-based advantage computation in GRPO is well-matched to the comparative paradigm on which many reward models are trained, i.e., discriminating between responses for the same prompt. It is also important to observe that GRPO employs a KL regularization term between the trained policy and the reference policy directly within the objective, as opposed to including it in the reward function, which keeps the calculation of $\hat{A}_{i,t}$ straightforward. Distinct from the KL penalty expression in (\ref{eq:PPO-reward}), here the KL divergence is estimated using an unbiased

\begin{algorithm}[t]
  \small
  \caption{Iterative Group Relative Policy Optimization}
  \textbf{Input} initial policy model $\pi_{\theta_{\text{init}}}$; reward models $r_\varphi$; task prompts $\mathcal{D}$; 
  hyperparameters $\varepsilon$, $\beta$, $\mu$
  \begin{algorithmic}[1]
    \State Initialize the policy model: $\pi_\theta \leftarrow \pi_{\theta_{\text{init}}}$
    \For{iteration = 1, \dots, I}
      \State Set the reference model: $\pi_{ref} \leftarrow \pi_{\theta}$
      \For{step = 1, \dots, M}
        \State Draw a mini-batch $\mathcal{D}_b$ from $\mathcal{D}$
        \State Store current policy: $\pi_{\theta_{old}} \leftarrow \pi_{\theta}$ 
        \State For each question $q \in \mathcal{D}_b$, sample $G$ responses $\{o_i\}_{i=1}^G$ according to $\pi_{\theta_{old}} (\cdot \mid q) $
        \State For each generated output $o_i$, compute its reward $r_i$ using $r_{\varphi}$
        \State Estimate the group-relative advantage $\hat{A}_{i,t}$ for the $t$-th token in $o_i$
        \For{GRPO iteration = 1, \dots, $\mu$}
          \State Update the policy $\pi_{\theta}$ by maximizing the GRPO loss (Equation \ref{eq:GC-GRPO})
        \EndFor
      \EndFor
      \State Refine $r_\varphi$ with continual training employing a replay buffer mechanism.
    \EndFor 
  \end{algorithmic}
  \textbf{Output} $\pi_\theta$
  \label{alg:iter-grpo}
\end{algorithm}

\subsubsection{Outcome Supervision RL with GRPO} Formally, for a given question $q$, a collection of $G$ outputs $\{o_1, o_2, \cdots, o_G\}$ is produced by sampling from the previous policy $\pi_{\theta_{old}}$. These outputs are evaluated by a reward model, resulting in reward assignments $\mathbf{r}=\{r_1, r_2, \cdots, r_G\}$. The rewards are standardized by centering and scaling, specifically by subtracting their mean and dividing by their standard deviation. With outcome supervision, the resulting normalized reward is provided only at the conclusion of each output $o_i$; all tokens within $o_i$ are assigned the same advantage, defined as $\hat{A}_{i, t} = \widetilde{r}_i = \frac{r_i- {\rm mean}(\mathbf{r})}{{\rm std}(\mathbf{r})}$. The policy is subsequently improved by optimizing the objective specified in equation (\ref{eq:GRPO-obj}).

\subsubsection{Process Supervision RL with GRPO} In contrast, outcome supervision yields feedback only at the sequence end, which may not adequately guide policies for complex reasoning required in mathematical problems. Motivated by \cite{wang2023math}, we further consider process supervision, which allocates rewards at each reasoning step. More precisely, for a question $q$ and $G$ samples $\{o_1, o_2, \cdots, o_G\}$, a dedicated process reward model evaluates every step, providing step-level rewards: $\mathbf{R} = \{ \{r_1^{index(1)},\cdots,r_1^{index(K_1)}\}, \cdots,  \{r_G^{index(1)},\cdots,r_G^{index(K_G)}\} \}$. Here, $index(j)$ denotes the final token index of step $j$, and $K_i$ represents the number of steps for output $o_i$. All rewards are normalized by computing $\widetilde{r}_i^{index(j)} = \frac{r_i^{index(j)} - {\rm mean(\mathbf{R})}}{{\rm std(\mathbf{R})}}$. Next, process supervision determines the token-level advantage by aggregating the normalized rewards from all subsequent steps: $\hat{A}_{i, t} = \sum_{index(j) \ge t} \widetilde{r}_i^{index(j)}$. The policy is then updated by maximizing the objective stated in equation (\ref{eq:GRPO-obj}).

\subsubsection{Iterative RL with GRPO} During the course of reinforcement learning, the existing reward model may become inadequate for effectively guiding the evolving policy model. To address this, we investigate iterative RL employing GRPO. As outlined in Algorithm \ref{alg:iter-grpo}, the iterative GRPO procedure involves producing fresh training datasets for the reward model using samples generated from the current policy model. The reward model undergoes continual updating with these new examples, while leveraging a replay buffer that preserves 10\% of the past data for stability. The reference model is regularly synchronized with the policy model, and subsequent training of the policy model is carried out using the updated reward model.

\subsection{Training and Evaluating DeepSeekMath-RL}

Our reinforcement learning experiments utilize DeepSeekMath-Instruct 7B as the base model. The RL training dataset consists exclusively of chain-of-thought questions sampled from GSM8K and MATH problems within the SFT corpus, totaling approximately 144K items. We omit other SFT questions in order to directly examine how RL affects benchmarks that are not represented in the RL training data. The construction of the reward model's training set follows the methodology from \citep{wang2023math}. The initial reward model is finetuned from DeepSeekMath-Base 7B, employing a learning rate of 2e-5. In GRPO training, the policy model uses a learning rate of 1e-6, and the KL regularization coefficient is fixed at 0.04. Each training question receives $64$ sampled completions. Generation outputs are limited to a maximum sequence length of 1024, and training is conducted in batches of 1024 samples. Policy model updates are restricted to a single step per exploration cycle. For evaluation, DeepSeekMath-RL 7B is tested on the same suite of benchmarks used by DeepSeekMath-Instruct 7B. Here, chain-of-thought GSM8K and MATH items are considered as in-domain, whereas all other benchmarks are classified as out-of-domain.

Table \ref{tab:sft_rl_math} presents a comparison of open- and closed-source models using both chain-of-thought and tool-augmented reasoning approaches across English and Chinese datasets. Our analysis reveals the following: (1) When leveraging chain-of-thought reasoning, DeepSeekMath-RL 7B achieves accuracy rates of 88.2\% on GSM8K and 51.7\% on MATH. These results exceed the performance of every open-source model from the 7B to 70B parameter range and surpass the majority of proprietary models as well. (2) Importantly, DeepSeekMath-RL 7B’s training was limited to chain-of-thought style instruction tuning data from GSM8K and MATH, and it was initialized from DeepSeekMath-Instruct 7B. Even with this narrow training regime, the model delivers higher scores than DeepSeekMath-Instruct 7B across all evaluated dimensions, underlining the impact and efficiency of reinforcement learning.

\section{Discussion}
In this section, we discuss the insights gained from our experiments involving both pre-training and reinforcement learning.

\subsection{Lessons Learnt in Pre-Training}

We begin by summarizing our observations regarding the pre-training phase. Unless noted otherwise, all experiments follow the methodology detailed in Section \ref{sec:quality-policy}. Additionally, throughout this discussion, the DeepSeekMath Corpus refers to an 89B-token collection obtained during the second data acquisition iteration.

\subsubsection{Code Training Benefits Mathematical Reasoning}

It is a widely held—though often untested—belief that training on code enhances reasoning abilities. Here, we provide empirical evidence specifically for the mathematics domain, showing that code-centric training bolsters a model’s proficiency in mathematical reasoning tasks, irrespective of the incorporation of tools.

To investigate the relationship between code training and mathematical reasoning, we explored both a two-stage training framework and a single-stage training approach:

\textbf{Two-Stage Training}

\begin{itemize}[topsep=0pt]
    \item \textbf{Code Training for 400B Tokens $\rightarrow$ Math Training for 150B Tokens}: DeepSeek-LLM 1.3B undergoes an initial phase of training on 400B code tokens, which is subsequently followed by a phase involving 150B math tokens.
    \item \textbf{General Training for 400B Tokens $\rightarrow$ Math Training for 150B Tokens}: For the baseline comparison, the first phase utilizes 400B general tokens (sampled from an extensive general corpus curated by DeepSeek-AI) rather than code tokens. This approach allows us to assess whether code-centric pretraining offers distinctive advantages over general pretraining for enhancing mathematical reasoning abilities.
\end{itemize}

\textbf{One-Stage Training}

\begin{itemize}[topsep=0pt]
    \item \textbf{Math Training for 150B Tokens}: The DeepSeek-LLM 1.3B model is directly trained for 150B math tokens, with no prior stage.
    \item \textbf{Training on a mixture of 400B Code Tokens and 150B Math Tokens}: Sequential math training after prior code training has been observed to impair coding competence. Therefore, we further explore whether jointly presenting code and math tokens in a single training phase not only augments mathematical reasoning but also counteracts catastrophic forgetting.
\end{itemize}

\paragraph{Results}
Table \ref{tab:code-math} and Table \ref{tab:code-math-general-eval} present the impact of various training strategies on downstream task performance.

Incorporating code training positively influences program-based mathematical reasoning under both multi-phase and single-phase regimens. As seen in Table \ref{tab:code-math}, the sequential approach—starting with code and transitioning to math—substantially boosts problem-solving performance on GSM8K and MATH when using Python, even prior to any explicit mathematical training. The addition of math-specific training in the subsequent stage further heightens performance. Notably, in the single-phase scenario, interleaving code and math tokens effectively reduces catastrophic forgetting often observed in the staged approach, while simultaneously enhancing coding (Table \ref{tab:code-math-general-eval}) and program-supported mathematical reasoning (Table \ref{tab:code-math}).

Moreover, code pretraining also advances mathematical reasoning in settings devoid of external tool use. During the staged training regime, the code pretraining stage alone achieves moderate improvements and increases the efficacy of the ensuing math training, culminating in optimal outcomes. In contrast, one-stage joint training with both code and math tokens is less advantageous for mathematical reasoning when external tools are absent. One possible explanation is that the parameter limitations of DeepSeek-LLM 1.3B may restrict its capacity to integrate both modalities simultaneously.

\subsubsection{Limited Effectiveness of ArXiv Papers in Enhancing Mathematical Reasoning}

Incorporating arXiv papers into mathematical pre-training datasets is a widespread practice \citep{minerva,gpt-f,llemma,mathpile}. Nevertheless, there is a lack of thorough investigation into their tangible contributions to models’ mathematical reasoning capabilities. Somewhat surprisingly, our experimental results indicate that arXiv papers may not substantially aid mathematical reasoning. Our study examines various model sizes, such as DeepSeek-LLM 1.3B and DeepSeek-Coder-Base-v1.5 7B \citep{deepseek-coder}, each trained on differently processed arXiv datasets:

\begin{itemize}[topsep=0pt]
    \item \textbf{MathPile} \citep{mathpile}: an 8.9B-token collection curated via rigorous cleaning and filtering, with over 85\% sourced from scientific arXiv documents;
    \item \textbf{ArXiv-RedPajama} \citep{redpajama}: comprising all arXiv LaTeX source files, systematically stripped of preambles, comments, macros, and bibliographies, amounting to 28.0B tokens in total.
\end{itemize}

We perform independent pre-training runs for DeepSeek-LLM 1.3B (150B tokens) and DeepSeek-Coder-Base-v1.5 7B (40B tokens) on each arXiv dataset. The outcomes consistently show that arXiv corpora, when used in isolation, do not meaningfully enhance— and may sometimes even impair— performance across a suite of mathematical evaluation benchmarks spanning varying complexity. These benchmarks include datasets for quantitative reasoning (GSM8K, MATH, see Table \ref{tab:arxiv-cot}), multiple-choice scientific questions (MMLU-STEM, also in Table \ref{tab:arxiv-cot}), as well as tasks from formal mathematics (miniF2F, Table \ref{tab:arxiv_atp}).

It is important, however, to acknowledge several caveats to these findings. Our investigation does not encompass:

\begin{itemize}[topsep=0pt]
    \item Potential influence of arXiv tokens on other mathematical tasks absent from our benchmarks, such as translating formal theorems or proofs into informal language;
    \item The impact of integrating arXiv data with other data sources in training;
    \item The possibility that scaling up model size could reveal benefits not seen in our present experiments.
\end{itemize}

Consequently, additional research is warranted to address these questions, which we defer for future work.

\subsection{Insights of Reinforcement Learning}
\subsubsection{Toward a Unified Framework} This section introduces a unified framework for interpreting a range of training approaches, including SFT, RFT, DPO, PPO, and GRPO. We also empirically assess how factors in this unified view influence training outcomes. More formally, for a given training algorithm, its gradient with respect to the model parameters $\theta$ is typically expressed as:

\begin{equation}
    \nabla_{\theta}\mathcal{J}_{\textcolor{red}{\mathcal{A}}}(\theta) = \mathbb{E}[\underbrace{(q,o) \sim \textcolor{red}{\mathcal{D}}}_{Data \ Source}]\left( \frac{1}{|o|} \sum_{t=1}^{|o|}  \underbrace{GC_{{\mathcal{A}}}(q, o, t, \textcolor{red}{\pi_{{rf}}})}_{Gradient \ Coefficient}  \nabla_{\theta}\log \pi_{\theta}(o_t | q, o_{<t})\right).
\label{eq:objective}
\end{equation}

This framework is composed of three principal elements: 1) \textit{Data Source $\mathcal{D}$}, indicating where the training samples are drawn from; 2) \textit{Reward Function $\pi_{{rf}}$}, which specifies how rewards are generated for learning; and 3) \textit{Algorithm $\mathcal{A}$}, responsible for converting data and reward information into the gradient coefficient $GC$, thereby adjusting the extent of reward or punishment for specific data. The following discussion organizes various prominent approaches under this generalized formulation:

\begin{itemize}
    \item  \textbf{Supervised Fine-tuning (SFT)}: This method involves updating a pretrained model using curated SFT data selected by human annotators.
    \item \textbf{Rejection Sampling Fine-tuning (RFT)}: RFT performs additional fine-tuning on the SFT model by selecting from its own generated outputs, retaining only those responses verified as correct for further training.
    \item \textbf{Direct Preference Optimization (DPO)}: DPO further trains the SFT model with pairs of augmented outputs generated from SFT, optimizing using the pairwise DPO loss to capture preferences.
    \item \textbf{Online Rejection Sampling Fine-tuning (Online RFT)}: In contrast to standard RFT, this strategy initializes the policy with the SFT model, but dynamically updates it by sampling and fine-tuning on outputs produced by the live policy itself.
    \item \textbf{PPO/GRPO}: Here, the SFT model serves as the starting point, and is strengthened via outputs sampled directly from the evolving policy model.
\end{itemize}

The elements comprising these methodologies are outlined in Table \ref{tab:data_policy}. For an in-depth derivation, readers are directed to Appendix \ref{app:analysis-rl}.

\paragraph{Data Source Analysis} The data sources are classified into online and offline sampling categories. Online sampling refers to utilizing training data generated dynamically through exploration by the current policy model, whereas offline sampling utilizes data derived from the initial SFT model’s output. RFT and DPO are representative of the offline paradigm, while Online RFT and GRPO exemplify the online approach.

According to the results illustrated in Figure \ref{fig:rl-analysis}, Online RFT demonstrates a marked performance improvement over RFT across both benchmarks. Initially, the performance of Online RFT closely matches that of RFT, owing to the similarity between the actor and the SFT model at the outset, with only slight distinctions in their sampled data. However, as training proceeds, the divergence in sampled data increases, and the advantages of real-time, online data collection become more pronounced—allowing Online RFT to substantially outperform RFT in the later stages.

\paragraph{Analysis of Gradient Coefficient} In these algorithms, the reward signal is processed into a gradient coefficient, which in turn influences updates to the model parameters. The reward function used is categorized in our studies as either `Rule'—assessing the quality of a response by verifying the correctness of the answer—or `Model'—where a trained reward model assigns scores to responses, with the training labels for the reward model determined by the rule-based approach. Notably, as reflected in Equations \ref{eq:GC-RFT} and \ref{eq:GC-GRPO}, a central distinction arises between GRPO and Online RFT: GRPO customizes its gradient coefficients based on the reward model’s outputs, permitting more nuanced reinforcement or penalization of responses in accordance with their reward magnitude. By contrast, Online RFT does not differentiate in this manner: it treats all correct responses equivalently, applying a uniform level of reinforcement, and does not penalize incorrect answers.

As illustrated in Figure \ref{fig:rl-analysis}, GRPO demonstrates superior performance over online RFT, emphasizing the effectiveness of adjusting positive and negative gradient coefficients. Moreover, the results show that GRPO+PS achieves better outcomes than GRPO+OS, underscoring the value of adopting more nuanced, step-specific gradient coefficients. The study also examines the use of iterative RL; in our experimental setup, we carry out two iterations. Figure \ref{fig:iter-rl} reveals that employing iterative RL yields notable performance gains, particularly following the initial iteration.

\subsubsection{Why RL Works?} In this work, reinforcement learning is applied using a subset of the instruction tuning dataset, yielding marked improvement over the base instruction tuning model. To deepen our understanding of RL's efficacy, we compare the Pass@K and Maj@K accuracies for both Instruct and RL models on two separate benchmarks. According to Figure \ref{fig:maj-pass}, RL boosts Maj@K results, but does not substantially affect Pass@K. This pattern suggests that RL fortifies the overall output distribution of the model, implying that \textbf{the observed performance boost mainly stems from elevating the likelihood of selecting the correct answer from among the TopK responses, rather than strengthening underlying model skills.} A similar \textbf{misalignment problem} was identified in reasoning tasks for SFT models by \citep{wang2023making}, with findings showing that preference alignment strategies \citep{yuan2023rrhf,song2023preference,wang2023making} can enhance their reasoning abilities.

\subsubsection{How to Achieve More Effective RL?}
\label{sec:effec_rl}
Our findings confirm that RL yields considerable improvements in mathematical reasoning problems. We additionally introduce a unified perspective for interpreting major training strategies, characterizing all as either explicit or simplified RL approaches within a single framework. According to Equation \ref{eq:objective}, three foundational elements are identified: Data Source, Algorithm, and Reward Function. Below, we outline some promising directions for each component.

\paragraph{Data Source} The data source serves as the cornerstone of any training procedure. Within the RL context, we define the data source as unlabeled questions paired with outputs sampled from the policy model. This paper exclusively employs prompts used in instruction tuning alongside simple nucleus sampling to generate outputs. We surmise that this choice partially explains why our RL framework mainly enhances Maj@K accuracy. Looking ahead, we plan to extend our RL pipeline to cover out-of-distribution questions and integrate \textbf{advanced sampling (decoding) strategies}, such as those utilizing tree-search algorithms \citep{yao2023tree}. In addition, \textbf{efficient inference techniques} \citep{xia-etal-2023-speculative,leviathan2023fast,kwon2023efficient,xia2024unlocking}—which influence how efficiently policy models can explore—are also expected to significantly affect outcomes.

\paragraph{Algorithms} Algorithms leverage both the data and reward signals to compute the gradients necessary for updating model parameters. Referring to Equation \ref{eq:objective}, existing approaches predominantly \textbf{TRUST} the feedback from the reward function, adjusting the conditional probability of generating specific tokens based on this signal. Nonetheless, absolute reliability of the reward signal cannot be guaranteed, especially for intricate tasks. For instance, even meticulously curated resources such as the PRM800K datasets \citep{lightman2023let}, annotated by experienced professionals, still exhibit roughly 20\% erroneous labels\footnote{\url{https://github.com/openai/prm800k/issues/12\#issuecomment-1728491852}}. Consequently, there is a need to investigate reinforcement learning algorithms that demonstrate resilience to imperfect or noisy reward information. We anticipate that alignment strategies inspired by \textbf{WEAK-TO-STRONG} \citep{burns2023weak} will fundamentally reshape algorithmic learning processes.

\paragraph{Reward Function} The reward function acts as the principal provider of learning signals during training. Within RL, this is typically represented by a neural reward model. We identify three pivotal avenues for advancing reward models: 1) \textbf{Strategies for improving the reward model's generalization capacity.} For RL to enhance the underlying abilities of LLMs—rather than merely reinforcing their existing output distributions—the reward model needs to generalize well to novel, out-of-distribution questions and complex outputs; 2) \textbf{Methods to encode and utilize the reward model's uncertainty.} Capturing model uncertainty could facilitate integration between suboptimal reward models and weak-to-strong learning frameworks; 3) \textbf{Development of efficient high-quality process reward models} capable of providing granular, process-level feedback, thereby supporting detailed training for reasoning-based tasks \citep{lightman2023let,wang2023math}.

\section{Conclusion, Limitation, and Future Work} 

In this work, we introduce DeepSeekMath, a model that surpasses all open-source alternatives on the competition-level MATH benchmark and closely rivals the performance of proprietary systems. DeepSeekMath~is initialized from DeepSeek-Coder-v1.5 7B and is continually trained over 500B tokens, with a substantial portion—120B math tokens—curated from Common Crawl. Our thorough ablation analysis demonstrates that web data serves as a highly promising resource for quality mathematical corpora, whereas the benefits derived from arXiv data appear more limited than anticipated. We propose Group Relative Policy Optimization (GRPO), a modified version of Proximal Policy Optimization (PPO), which offers notable enhancements in mathematical reasoning efficiency alongside reduced memory overhead. Experimental findings indicate that GRPO can further improve the performance of DeepSeekMath-Instruct 7B, even when it already performs strongly on standard benchmarks. Additionally, we introduce a unifying perspective to relate existing reinforcement learning strategies and outline several promising avenues for advancing reinforcement learning methodologies.

Despite DeepSeekMath’s~strong quantitative reasoning abilities, its proficiency in areas such as geometry and formal proof lags behind that of closed-source models. For example, our dry-run analysis revealed the model’s struggles with triangle and ellipse problems, suggesting possible data selection biases during pre-training and fine-tuning. Furthermore, due to its model scale, DeepSeekMath~is outperformed by GPT-4 in few-shot learning scenarios; while GPT-4 shows marked improvements with few-shot prompts, DeepSeekMath~exhibits comparable results in both zero-shot and few-shot settings. Future efforts will concentrate on refining our engineered data selection workflow to assemble higher-quality pre-training datasets. Moreover, we plan to investigate the promising directions discussed in Section \ref{sec:effec_rl} for developing more robust reinforcement learning approaches tailored to LLMs.

\end{document}